\documentclass[11pt]{article}

\usepackage[T1]{fontenc}
\usepackage[utf8]{inputenc}
\usepackage{amsmath,amssymb,amsthm,mathtools}
\usepackage{microtype}
\usepackage{hyperref}
\hypersetup{hidelinks}
\usepackage{geometry}
\geometry{margin=1in}

\newtheorem{theorem}{Theorem}[section]
\newtheorem{proposition}[theorem]{Proposition}
\newtheorem{lemma}[theorem]{Lemma}
\newtheorem{corollary}[theorem]{Corollary}
\newtheorem{remark}[theorem]{Remark}
\newtheorem{definition}[theorem]{Definition}

\newcommand{\C}{\mathbb{C}}
\newcommand{\op}{\mathrm{op}}
\newcommand{\cmax}{C_{\max}}
\newcommand{\gmax}{G_{\max}}
\newcommand{\dd}{\delta}
\newcommand{\etaq}{\eta}
\newcommand{\tauq}{\tau}
\newcommand{\PhiRho}{\Phi_{\rho}}

\title{Rigidity and Scale-Sharp Stability for Gaps of\\
Finite Orthogonal Decompositions}

\author{Bertrand Dalimier\\
Independent Researcher}

\date{}

\begin{document}

\maketitle

\begin{abstract}
Let $(P_i)_{i\in I}$ and $(Q_i)_{i\in I}$ be two labelled finite orthogonal
resolutions of the identity on a finite-dimensional complex Hilbert space,
with $\operatorname{rank}P_i=\operatorname{rank}Q_i$ for every $i$.  Put
\[
  p_i=\|P_i-Q_i\|_{\op}^2,\qquad
  \dd^2=\sum_i p_i,
\]
and, for $S\subseteq I$,
\[
  P_S=\sum_{i\in S}P_i,\qquad
  Q_S=\sum_{i\in S}Q_i,\qquad
  \gmax=\max_{S\subseteq I}\|P_S-Q_S\|_{\op}.
\]
We prove the envelope $\gmax^2\le\dd^2/2$ and characterize positive-defect
equality by the fact that exactly two cell gaps are nonzero.  The main result
is a cardinality- and dimension-free inverse theorem: if
$0\le\eta\le1/5$, $\dd>0$, and
$\gmax^2\ge(1-\eta)\dd^2/2$, then the optimal relative defect left after
retaining two cells satisfies
\[
  \tau\le
  2\eta+\frac{36\eta}{(1-3\eta)\dd^2}.
\]
Thus the theorem is uniform on every positive defect scale.  A one-parameter
three-cell family satisfies
\[
  \dd_t^2\to0,\qquad \eta_t\to0,\qquad
  \dd_t^2\,\frac{\tau_t}{\eta_t}\to4,
\]
showing that the inverse-square dependence on the defect scale is necessary
up to constants.  Equivalently, the problem concerns unitary conjugation:
if $Q_i=UP_iU^*$, then
$\|P_i-Q_i\|_{\op}=\|[P_i,U]\|_{\op}$.  The unitary theorem chain and a
discrete specialization of the scale obstruction are machine-checked in
Lean 4.
\end{abstract}

\noindent\textbf{Keywords.}
orthogonal projections; gap between subspaces; principal angles;
near-equality stability; rigidity; commutator inequalities; Lean 4

\noindent\textbf{2020 Mathematics Subject Classification.}
15A60 (primary); 15A45, 47A63 (secondary).

\section{Introduction}

Let
\[
  \mathcal P=(P_i)_{i\in I},
  \qquad
  \mathcal Q=(Q_i)_{i\in I}
\]
be two labelled finite orthogonal decompositions of a finite-dimensional
complex Hilbert space.  Assume that corresponding summands have equal
dimension.  The operator norm
\[
  \|P_i-Q_i\|_{\op}
\]
is the standard gap between the two corresponding subspaces; for
equal-dimensional subspaces it is the sine of their largest principal angle
\cite{BjorckGolub1973,MiaoBenIsrael1992,GalantaiHegedus2006,Galantai2008}.
This makes the family
\[
  p_i:=\|P_i-Q_i\|_{\op}^2
\]
a natural quadratic profile of cellwise subspace displacement.

The equal-rank hypothesis is equivalent to the existence of a unitary
$U$ such that
\[
  Q_i=UP_iU^*
  \qquad(i\in I).
\]
For any such implementing unitary,
\begin{equation}\label{eq:gap-comm}
  [P_i,U]U^*=P_i-UP_iU^*=P_i-Q_i,
  \qquad
  \|[P_i,U]\|_{\op}=\|P_i-Q_i\|_{\op}.
\end{equation}
Thus the same problem can be read either as geometry of two orthogonal
decompositions or as a finite profile of projector--unitary commutators.

For a set of labels $S\subseteq I$, define the corresponding coarsenings
\[
  P_S:=\sum_{i\in S}P_i,
  \qquad
  Q_S:=\sum_{i\in S}Q_i,
\]
and their gap
\[
  G_S:=\|P_S-Q_S\|_{\op}.
\]
We study the maximal coarsening gap
\[
  \gmax:=\max_{S\subseteq I}G_S
\]
relative to the global quadratic cell budget
\[
  \dd^2:=\sum_{i\in I}p_i.
\]
The elementary envelope is
\begin{equation}\label{eq:envelope}
  \gmax^2\le\frac{\dd^2}{2}.
\end{equation}
The principal contribution of the paper is not the scalar factor $1/2$ in
\eqref{eq:envelope}, but the structure of its extremizers and
near-extremizers.

There are three results.  First, when $\dd>0$, equality in
\eqref{eq:envelope} holds exactly when two cells, and only two cells, have
nonzero gap.  Second, near equality forces an approximate version of the same
two-cell concentration.  If
\[
  \gmax^2\ge(1-\eta)\frac{\dd^2}{2},
  \qquad 0\le\eta\le\frac15,
\]
then the relative defect outside the best two cells is at most
\[
  2\eta+\frac{36\eta}{(1-3\eta)\dd^2}.
\]
The constants are independent of the dimension of the Hilbert space and of
the number of cells.  Third, a one-parameter rank-one example in dimension
three shows that the factor $\dd^{-2}$ is unavoidable up to multiplicative
constants: along the example,
\[
  \dd_t^2\,\frac{\tau_t}{\eta_t}\longrightarrow4.
\]

The proof separates into two ingredients.  A cut estimate compares a
coarsening gap simultaneously with the quadratic budgets on the two sides of
the cut.  Near equality then produces normalized witness vectors whose Gram
matrix is close to that of an orthonormal family.  A cardinality-free Gram
estimate forces one dominant cell on each side of the maximizing cut.

The gap formulation places the result in the classical geometry of pairs of
subspaces and orthogonal projections
\cite{Halmos1969TwoSubspaces,BottcherSpitkovsky2010,Galantai2008}.
The relation to commutator geometry, almost-commuting problems, and recent
work on pairs of projections is discussed in Section~\ref{sec:related}.

The unitary theorem chain has a companion Lean 4 formalization.  Formal
verification is described briefly in Section~\ref{sec:formal}; it is used as
a verification artifact rather than as a premise of the mathematical
argument.

% TODO before submission:
% 1. Replace repository commit by archival release/DOI.
% 2. Apply final LAA/Elsevier template if required.
% 3. Prepare cover letter and submission metadata.

\section{Two labelled orthogonal decompositions}\label{sec:setup}

Let $H$ be a finite-dimensional complex Hilbert space and let $I$ be a finite
index set.  Let $(P_i)_{i\in I}$ and $(Q_i)_{i\in I}$ be orthogonal
projections satisfying
\begin{equation}\label{eq:decomp}
  P_iP_j=0,\quad Q_iQ_j=0\quad(i\ne j),
  \qquad
  \sum_{i\in I}P_i=\sum_{i\in I}Q_i=I_H,
\end{equation}
and assume
\[
  \operatorname{rank}P_i=\operatorname{rank}Q_i
  \qquad(i\in I).
\]
Choose, when convenient, a unitary $U$ satisfying $Q_i=UP_iU^*$ for every
$i$; such a unitary exists by choosing unitary identifications between the
corresponding summands and taking their orthogonal direct sum.

Define
\begin{equation}\label{eq:pi}
  p_i:=\|P_i-Q_i\|_{\op}^2
      =\|[P_i,U]\|_{\op}^2
\end{equation}
and
\begin{equation}\label{eq:delta}
  \dd^2:=\sum_{i\in I}p_i.
\end{equation}
For $S\subseteq I$, set
\begin{equation}\label{eq:cut}
  P_S:=\sum_{i\in S}P_i,\qquad
  Q_S:=\sum_{i\in S}Q_i,\qquad
  G_S:=\|P_S-Q_S\|_{\op}.
\end{equation}
By \eqref{eq:gap-comm},
\[
  G_S=\|[P_S,U]\|_{\op}.
\]
Finally,
\[
  \gmax:=\max_{S\subseteq I}G_S.
\]

If $\dd>0$, define the optimal relative two-cell tail
\begin{equation}\label{eq:tau}
  \tau:=
  \min_{\substack{i,j\in I\\i\ne j}}
  \frac{\dd^2-p_i-p_j}{\dd^2}.
\end{equation}
Thus $\tau$ is the fraction of the global quadratic gap budget left after the
best pair of labelled cells is retained.

The quantities in this formulation depend only on the pair
$(\mathcal P,\mathcal Q)$, not on the choice of implementing unitary $U$.
The unitary is used only as a convenient proof device.

\section{The cut envelope}\label{sec:envelope}

We first isolate the elementary part of the argument.  In fact the cut
estimate itself does not require unitarity.

\begin{proposition}[Operator cut estimate]\label{prop:operator-envelope}
Let $(P_i)_{i\in I}$ be a finite orthogonal resolution of the identity and
let $T$ be any linear operator on $H$.  Put
\[
  d_i:=\|[P_i,T]\|_{\op},
  \qquad
  D^2:=\sum_{i\in I}d_i^2.
\]
For $S\subseteq I$ and $P_S=\sum_{i\in S}P_i$,
\begin{equation}\label{eq:cut-min}
  \|[P_S,T]\|_{\op}^2
  \le
  \min\left\{
    \sum_{i\in S}d_i^2,\,
    \sum_{i\notin S}d_i^2
  \right\}.
\end{equation}
Consequently,
\[
  \max_{S\subseteq I}\|[P_S,T]\|_{\op}^2
  \le\frac{D^2}{2}.
\]
\end{proposition}

\begin{proof}
Write $P=P_S$ and $P^\perp=I-P$.  Relative to
$H=PH\oplus P^\perp H$,
\[
  [P,T]
  =
  \begin{pmatrix}
    0 & PTP^\perp\\
    -P^\perp TP & 0
  \end{pmatrix},
\]
so its norm is the maximum of the norms of the two cross blocks.

We first control both blocks by the budget on $S$.  If
$x\in P^\perp H$, then $P_i x=0$ for $i\in S$, hence
$P_iTx=[P_i,T]x$.  Orthogonality of the ranges $P_iH$ gives
\[
  \|PTP^\perp x\|^2
  =
  \sum_{i\in S}\|P_iTx\|^2
  \le
  \left(\sum_{i\in S}d_i^2\right)\|x\|^2.
\]
If instead $x\in PH$, write $x=\sum_{i\in S}x_i$ with
$x_i=P_i x$.  Since
$P^\perp TP_i=-P^\perp[P_i,T]P_i$,
\[
  \|P^\perp TPx\|
  \le
  \sum_{i\in S}d_i\|x_i\|
  \le
  \left(\sum_{i\in S}d_i^2\right)^{1/2}\|x\|
\]
by Cauchy--Schwarz.  Thus
$\|[P,T]\|_{\op}^2\le\sum_{i\in S}d_i^2$.
Applying the same argument to the complementary resolution gives the bound by
$\sum_{i\notin S}d_i^2$, proving \eqref{eq:cut-min}.  The final estimate is
the scalar inequality $\min\{x,D^2-x\}\le D^2/2$.
\end{proof}

Applying Proposition~\ref{prop:operator-envelope} to an implementing unitary
and using \eqref{eq:gap-comm} gives the geometric form.

\begin{theorem}[Gap envelope]\label{thm:envelope}
For the two labelled decompositions of Section~\ref{sec:setup},
\begin{equation}\label{eq:gap-envelope}
  \boxed{\gmax^2\le\frac{\dd^2}{2}}.
\end{equation}
The coefficient $1/2$ is attained.
\end{theorem}

The rest of the paper is concerned with the inverse question: what must the
cellwise gap profile look like when \eqref{eq:gap-envelope} is nearly an
equality?

\section{Quantitative stability of near-saturators}\label{sec:stability}

Fix an implementing unitary $U$ with $Q_i=UP_iU^*$.  All quantities in the
statement remain the intrinsic gap quantities of Section~\ref{sec:setup}.

For $S\subseteq I$, write
\[
  A_S:=\sum_{i\in S}p_i,
  \qquad
  G:=G_S.
\]

\subsection{A cardinality-free Gram lemma}

The following elementary lemma is the concentration mechanism used below.

\begin{lemma}[Gram spread]\label{lem:gram}
Let $(\xi_i)_{i\in J}$ be an orthonormal finite family in a complex Hilbert
space, let $(\zeta_i)_{i\in J}$ be any family, and let $p_i\ge0$.  Suppose
that for $i\ne j$,
\[
  p_ip_j=|\langle\zeta_i,\zeta_j\rangle|^2.
\]
Set
\[
  q:=\sum_{i\in J}\|\zeta_i-\xi_i\|^2.
\]
Then
\begin{equation}\label{eq:gram-spread-general}
  \sum_{i\in J}\sum_{\substack{j\in J\\j\ne i}}p_ip_j
  \le 6q+3q^2.
\end{equation}
In particular, if $0\le e\le1/2$ and $q\le2e$, then the left-hand side is
at most $18e$.
\end{lemma}

\begin{proof}
Put $e_i:=\zeta_i-\xi_i$.  For $i\ne j$, orthogonality gives
\[
  \langle\zeta_i,\zeta_j\rangle
  =
  \langle e_i,\xi_j\rangle
  +\langle\xi_i,e_j\rangle
  +\langle e_i,e_j\rangle.
\]
Hence $(a+b+c)^2\le3(a^2+b^2+c^2)$ implies
\[
  p_ip_j
  \le
  3\left(
    |\langle e_i,\xi_j\rangle|^2+
    |\langle\xi_i,e_j\rangle|^2+
    |\langle e_i,e_j\rangle|^2
  \right).
\]
After summing over ordered pairs $i\ne j$, the first two double sums are each
at most $q$ by Bessel's inequality.  The third is at most
\[
  \sum_{i,j}\|e_i\|^2\|e_j\|^2=q^2
\]
by Cauchy--Schwarz.  This proves \eqref{eq:gram-spread-general}.  If
$q\le2e$ and $e\le1/2$, then
\[
  6q+3q^2\le12e+12e^2\le18e.
\]
\end{proof}

\subsection{One-sided concentration}

\begin{lemma}[One-sided concentration]\label{lem:onesided}
Let $S$ be a cut with $G>0$, and put
\[
  e:=A_S-G^2.
\]
If $e\le G^2/2$, then there exists $i\in S$ such that
\begin{equation}\label{eq:onesided}
  A_S-p_i
  \le
  2e+\frac{18e}{G^2-e}.
\end{equation}
\end{lemma}

\begin{proof}
Relative to $P_SH\oplus P_{S^c}H$, the commutator $[P_S,U]$ has two cross
blocks.  We treat the case in which the incoming block
$P_SU P_{S^c}$ has norm $G$; the other case is identical after exchanging
the two sides of the cut.

Choose a unit vector $v\in P_{S^c}H$ such that
$\|P_SUv\|=G$.  For $c\in S$, define
\[
  z_c:=U^*P_cUv,\qquad
  r_c:=\|z_c\|^2,\qquad
  \alpha_c:=p_c-r_c.
\]
The vectors $z_c$ are mutually orthogonal.  Moreover
$r_c\le p_c$, because on the complementary witness $v$ the vector
$P_cUv$ is the value of the cell commutator, and
\[
  \sum_{c\in S}r_c=\|P_SUv\|^2=G^2.
\]
Consequently
\begin{equation}\label{eq:witness-defect}
  \alpha_c\ge0,
  \qquad
  \sum_{c\in S}\alpha_c=e.
\end{equation}

Call $c$ bad if $\alpha_c>p_c/2$, and good otherwise.  From
\eqref{eq:witness-defect},
\begin{equation}\label{eq:badmass}
  \sum_{\text{bad }c}p_c\le2e.
\end{equation}
Let $J$ be the active good cells, so $p_c>0$ and
$r_c=p_c-\alpha_c\ge p_c/2>0$ on $J$.

For $c\in J$, put
\[
  q_c:=(I-P_c)z_c,\qquad
  \xi_c:=\frac{z_c}{\sqrt{r_c}},
  \qquad
  \zeta_c:=
  \frac{P_cz_c}{\sqrt{r_c}}+\sqrt{p_c}\,v.
\]
The family $(\xi_c)_{c\in J}$ is orthonormal.  Since
$v\in P_{S^c}H$, the vectors $P_cz_c$ lie in mutually orthogonal cells and
are orthogonal to $v$.  Therefore, for $c\ne d$,
\begin{equation}\label{eq:zeta-gram}
  |\langle\zeta_c,\zeta_d\rangle|^2=p_cp_d.
\end{equation}

We next estimate the distance from $\zeta_c$ to $\xi_c$.  The residual
$q_c$ satisfies
\begin{equation}\label{eq:residual-bounds}
  \|q_c\|^2\le p_cr_c,
  \qquad
  \operatorname{Re}\langle v,q_c\rangle=r_c.
\end{equation}
Indeed, $Uq_c=[P_c,U]z_c$, which gives the first inequality, while
$P_cv=0$ and
\[
  \langle v,z_c\rangle
  =
  \langle Uv,P_cUv\rangle
  =
  \|P_cUv\|^2=r_c
\]
give the second.  Since
\[
  \xi_c-\zeta_c
  =
  \frac{q_c}{\sqrt{r_c}}-\sqrt{p_c}\,v,
\]
\eqref{eq:residual-bounds} yields
\[
  \|\xi_c-\zeta_c\|^2
  \le
  2p_c-2\sqrt{p_cr_c}
  \le
  2(p_c-r_c)
  =
  2\alpha_c.
\]
Thus
\begin{equation}\label{eq:gram-close}
  \sum_{c\in J}\|\zeta_c-\xi_c\|^2\le2e.
\end{equation}

Because $G\le1$ for a difference of orthogonal projections and
$e\le G^2/2$, we have $e\le1/2$.  Lemma~\ref{lem:gram},
\eqref{eq:zeta-gram}, and \eqref{eq:gram-close} therefore give
\begin{equation}\label{eq:gram18}
  \sum_{c\in J}\sum_{\substack{d\in J\\d\ne c}}p_cp_d
  \le18e.
\end{equation}
The active-good mass
\[
  M:=\sum_{c\in J}p_c
\]
satisfies, by \eqref{eq:badmass},
\[
  M\ge A_S-2e=G^2-e=:L>0.
\]
Choose $i\in J$ with $p_i=\max_{c\in J}p_c$.  Since
\[
  \sum_{c\ne d}p_cp_d
  =
  M^2-\sum_cp_c^2
  \ge M^2-p_iM
  =
  M(M-p_i),
\]
\eqref{eq:gram18} implies
\[
  M-p_i\le\frac{18e}{M}\le\frac{18e}{L}.
\]
Adding the bad-cell mass \eqref{eq:badmass} proves
\eqref{eq:onesided}.
\end{proof}

\subsection{Bilateral concentration and the intrinsic stability bound}

For a cut $S$, write
\[
  A:=\dd^2,\qquad G:=G_S,\qquad
  e_S:=A_S-G^2,\qquad
  e_T:=A_{S^c}-G^2,
\]
and
\[
  E:=A-2G^2.
\]
Since $A_S+A_{S^c}=A$,
\begin{equation}\label{eq:e-split}
  e_S+e_T=E.
\end{equation}

\begin{proposition}[Bilateral cut concentration]\label{prop:bilateral}
Assume $G>0$ and $E\le G^2/2$.  Then there exist
$i\in S$ and $j\in S^c$ such that
\begin{equation}\label{eq:bilateral}
  A-p_i-p_j
  \le
  2E+\frac{18E}{G^2-E}.
\end{equation}
\end{proposition}

\begin{proof}
The cut estimate gives $G^2\le A_S$ and $G^2\le A_{S^c}$, hence
$e_S,e_T\ge0$.  By \eqref{eq:e-split}, both are at most $E$, so
Lemma~\ref{lem:onesided} applies on both sides of the cut.  We obtain
$i\in S$ and $j\in S^c$ with
\[
  A_S-p_i
  \le
  2e_S+\frac{18e_S}{G^2-e_S},
  \qquad
  A_{S^c}-p_j
  \le
  2e_T+\frac{18e_T}{G^2-e_T}.
\]
Let $d:=G^2-E>0$.  Since
$d\le G^2-e_S$ and $d\le G^2-e_T$,
\[
  \frac{e_S}{G^2-e_S}
  +
  \frac{e_T}{G^2-e_T}
  \le
  \frac{E}{d}.
\]
Adding the two inequalities proves \eqref{eq:bilateral}.
\end{proof}

\begin{theorem}[Quantitative two-cell stability]\label{thm:stability}
Assume $\dd>0$, $0\le\eta\le1/5$, and
\begin{equation}\label{eq:near}
  \gmax^2\ge(1-\eta)\frac{\dd^2}{2}.
\end{equation}
Then
\begin{equation}\label{eq:stability-intrinsic}
  \boxed{
  \tau\le
  2\eta+\frac{36\eta}{(1-3\eta)\dd^2}.
  }
\end{equation}
Equivalently, there exist distinct $i,j$ such that
\begin{equation}\label{eq:absolute-tail}
  \dd^2-p_i-p_j
  \le
  2\eta\dd^2+\frac{36\eta}{1-3\eta}.
\end{equation}
\end{theorem}

\begin{proof}
Choose a cut $S$ attaining $\gmax$ and put
\[
  A:=\dd^2,\qquad G:=\gmax,\qquad E:=A-2G^2.
\]
Theorem~\ref{thm:envelope} gives $E\ge0$, while
\eqref{eq:near} gives
\begin{equation}\label{eq:Eeta}
  E\le\eta A.
\end{equation}
Since $\eta\le1/5$ and $A=2G^2+E$,
\[
  E\le\frac A5
  \quad\Longrightarrow\quad
  E\le\frac{G^2}{2}.
\]
Also $G>0$.  Proposition~\ref{prop:bilateral} therefore gives distinct
$i,j$ with
\[
  A-p_i-p_j
  \le
  2E+\frac{18E}{G^2-E}.
\]
Because $G^2=(A-E)/2$,
\[
  G^2-E=\frac{A-3E}{2}
  \ge\frac{(1-3\eta)A}{2}.
\]
Using \eqref{eq:Eeta},
\[
  A-p_i-p_j
  \le
  2\eta A+\frac{36\eta}{1-3\eta},
\]
which is \eqref{eq:absolute-tail}.  Division by $A=\dd^2$ gives
\eqref{eq:stability-intrinsic}.
\end{proof}

\begin{corollary}[Uniform positive-scale form]\label{cor:rho}
If, in addition, $\dd\ge\rho>0$, then
\begin{equation}\label{eq:rho}
  \tau\le
  2\eta+\frac{36\eta}{(1-3\eta)\rho^2}.
\end{equation}
If also $\eta\le1/6$, then
\begin{equation}\label{eq:linear}
  \tau\le
  \eta\left(2+\frac{72}{\rho^2}\right).
\end{equation}
\end{corollary}

\begin{proof}
Use $\dd^{-2}\le\rho^{-2}$ in
\eqref{eq:stability-intrinsic}.  For $\eta\le1/6$,
$1-3\eta\ge1/2$.
\end{proof}

\section{Exact rigidity at saturation}\label{sec:rigidity}

Call a cell active when $p_i>0$.

\begin{theorem}[Exact equality classification]\label{thm:rigidity}
Assume $\dd>0$.  Then
\[
  \boxed{
  \gmax^2=\frac{\dd^2}{2}
  \quad\Longleftrightarrow\quad
  \text{exactly two cells are active}.
  }
\]
In the equality case, if $i$ and $j$ are the two active cells, then
\[
  p_i=p_j=\frac{\dd^2}{2}.
\]
\end{theorem}

\begin{proof}
Assume first that $\gmax^2=\dd^2/2$.  Apply
Theorem~\ref{thm:stability} with $\eta=0$.  Then $\tau=0$, so there are
distinct $i,j$ with
\[
  \dd^2-p_i-p_j=0.
\]
All other $p_k$ therefore vanish.  Since singleton sets are admissible cuts,
\[
  p_i\le\gmax^2=\frac{\dd^2}{2},
  \qquad
  p_j\le\frac{\dd^2}{2}.
\]
Together with $p_i+p_j=\dd^2$, this gives
$p_i=p_j=\dd^2/2>0$.

Conversely, suppose exactly two cells $i\ne j$ are active.  Then
$\dd^2=p_i+p_j$, so
\[
  \max\{p_i,p_j\}\ge\frac{\dd^2}{2}.
\]
A singleton cut gives
$\gmax^2\ge\max\{p_i,p_j\}$, while
Theorem~\ref{thm:envelope} gives the reverse bound
$\gmax^2\le\dd^2/2$.  Equality follows.
\end{proof}

The hypothesis $\dd>0$ only excludes the degenerate case
$\mathcal P=\mathcal Q$.

\section{A one-parameter three-cell scale obstruction}\label{sec:sharp}

We now show that the inverse-square factor in
Theorem~\ref{thm:stability} is not an artifact of the proof.  The example is
a one-parameter family of rank-one decompositions in $\C^3$.

Let $0<t\le1/2$ and set
\[
  A_t:=\sqrt{1-t},
  \qquad
  B_t:=\sqrt t,
\]
\[
  c_t:=1-t,
  \qquad
  s_t:=\sqrt{2t-t^2}.
\]
Thus $A_t^2+B_t^2=1$ and $c_t^2+s_t^2=1$.  Define the real orthogonal
matrices
\[
  W_t=
  \begin{pmatrix}
    A_t&-B_t&0\\
    B_t&A_t&0\\
    0&0&1
  \end{pmatrix},
  \qquad
  R_t=
  \begin{pmatrix}
    c_t&0&-s_t\\
    0&1&0\\
    s_t&0&c_t
  \end{pmatrix},
\]
and the unitary
\[
  U_t:=W_tR_tW_t^*.
\]
Let $P_0,P_1,P_2$ be the coordinate rank-one projections and put
\[
  Q_k(t):=U_tP_kU_t^*
  \qquad(k=0,1,2).
\]

For a unit vector $e$ and its rank-one projector $P_e$,
\begin{equation}\label{eq:rankone}
  \|P_e-UP_eU^*\|_{\op}^2
  =
  \|[P_e,U]\|_{\op}^2
  =
  1-|\langle e,Ue\rangle|^2.
\end{equation}
The diagonal entries of $U_t$ are
\[
  1-t(1-t),\qquad 1-t^2,\qquad 1-t.
\]
Hence the three squared cell gaps are
\begin{align}
  p_0(t)
  &=2t(1-t)-t^2(1-t)^2, \label{eq:p0}\\
  p_1(t)
  &=2t^2-t^4, \label{eq:p1}\\
  p_2(t)
  &=2t-t^2. \label{eq:p2}
\end{align}
Equivalently, with $d=t$ and $x\in[0,1]$, the first two expressions have
the form
\[
  f_d(x):=2dx-d^2x^2
\]
at $x=1-t$ and $x=t$, while $p_2=f_d(1)$.  Since $f_d$ is increasing on
$[0,1]$ for $0<d\le1$, and $t\le1-t$, we have
\begin{equation}\label{eq:ordering}
  p_1(t)\le p_0(t)\le p_2(t).
\end{equation}

For three cells every nontrivial cut is a singleton or the complement of a
singleton, and
\[
  P_{S^c}-Q_{S^c}=-(P_S-Q_S).
\]
Therefore
\begin{equation}\label{eq:gmax-sharp}
  \gmax(t)^2=p_2(t)=2t-t^2.
\end{equation}
Summing \eqref{eq:p0}--\eqref{eq:p2} gives
\begin{equation}\label{eq:global-sharp}
  \dd_t^2
  =
  4t-2t^2+2t^3(1-t).
\end{equation}
Define the relative saturation deficit by
\[
  \eta_t:=
  \frac{\dd_t^2-2\gmax(t)^2}{\dd_t^2}.
\]
Using \eqref{eq:gmax-sharp} and \eqref{eq:global-sharp},
\begin{equation}\label{eq:eta-sharp}
  \eta_t=
  \frac{2t^3(1-t)}{\dd_t^2}.
\end{equation}
By \eqref{eq:ordering}, the optimal two-cell tail omits cell $1$, so
\begin{equation}\label{eq:tau-sharp}
  \tau_t=
  \frac{2t^2-t^4}{\dd_t^2}.
\end{equation}
Consequently,
\begin{equation}\label{eq:ratio-sharp}
  \frac{\tau_t}{\eta_t}
  =
  \frac{2-t^2}{2t(1-t)}.
\end{equation}

\begin{theorem}[Scale-sharp obstruction]\label{thm:scale-sharp}
As $t\downarrow0$,
\begin{equation}\label{eq:scale-limit}
  \dd_t^2\longrightarrow0,
  \qquad
  \eta_t\longrightarrow0,
  \qquad
  \dd_t^2\,\frac{\tau_t}{\eta_t}\longrightarrow4.
\end{equation}
Hence for every $c<4$ and every $\varepsilon>0$ there exists
$0<t\le1/2$ such that
\[
  0<\dd_t^2<\varepsilon,
  \qquad
  0<\eta_t<\varepsilon,
  \qquad
  \dd_t^2\frac{\tau_t}{\eta_t}>c.
\]
In particular, no estimate $\tau\le K\eta$ can hold uniformly with $K$
independent of the defect scale, and no uniform estimate
$\tau\le\eta\,g(\dd)$ can have $g(\dd)=o(\dd^{-2})$ as $\dd\downarrow0$.
\end{theorem}

\begin{proof}
From \eqref{eq:global-sharp},
\[
  \frac{\dd_t^2}{t}\longrightarrow4.
\]
Equation \eqref{eq:eta-sharp} then gives $\eta_t\to0$, while
\eqref{eq:ratio-sharp} gives
\[
  t\,\frac{\tau_t}{\eta_t}
  =
  \frac{2-t^2}{2(1-t)}
  \longrightarrow1.
\]
Multiplying the last two limits gives the constant $4$ in
\eqref{eq:scale-limit}.  The quantified statement follows by choosing $t$
sufficiently small.  The two impossibility statements are immediate.
\end{proof}

Thus Theorem~\ref{thm:stability} has the correct inverse-square dependence
on the intrinsic defect $\dd$, up to the value of the universal
multiplicative constant.

\section{Formal verification}\label{sec:formal}

The unitary form of the envelope, the exact rigidity theorem, the
cardinality-free stability argument, the rank-one identity, and a discrete
three-cell specialization of the scale obstruction have been machine-checked
in Lean 4.  The public repository is
\begin{center}
\url{https://github.com/Bobart0/everettian-decoherence-lean}.
\end{center}
The formal development uses canonical coordinate spaces $\C^n$; the
coordinate-free finite-dimensional statements used here are obtained by
unitary identification.

The theorem-to-code map and a dedicated audit module are included in the
repository.  The formalization head preceding the manuscript branch is
\begin{center}
\small\texttt{81eb3ff53f6219749b84491513ff437fe7c545c1}.
\end{center}
Its continuous-integration build, audit aggregate, guards, and
diff-hygiene checks pass.

Two presentation choices in the present manuscript go beyond the literal
surface syntax of the formal artifact.  First, Section~\ref{sec:envelope}
states the elementary cut estimate for an arbitrary operator $T$, whereas
the publication-facing Lean theorem used downstream is its unitary
specialization.  Second, Section~\ref{sec:sharp} presents the sharpness
mechanism as the simpler one-parameter real curve $t\downarrow0$; the Lean
artifact verifies a discrete two-parameter specialization with the same
exact budget identities and the same limiting constant $4$.

\paragraph{AI-assisted formalization workflow.}
OpenAI ChatGPT was used as an interactive assistant during parts of the Lean
development, including proof-structure exploration, code drafting, and
diagnosis of compiler errors.  Every Lean declaration referred to as
machine-checked is checked by the Lean kernel and rebuilt by continuous
integration.  The author reviewed the mathematical statements and proof
dependencies and remains responsible for the correctness and interpretation
of the results.

\section{Discussion}

Theorem~\ref{thm:stability} is intrinsic in the observed gap scale $\dd$.
The auxiliary lower bound $\dd\ge\rho$ is therefore not part of the principal
statement; it is only a convenient way to obtain a modulus uniform on a
positive scale.  Corollary~\ref{cor:rho} recovers that uniform formulation.

The theorem is dimension-free in two senses: its constants do not depend on
$\dim H$ and do not depend on $|I|$.  The price is the factor
$\dd^{-2}$ as the two decompositions coalesce.  Section~\ref{sec:sharp}
shows that this loss is structural.  The constants $36$ and $72$ are not
claimed to be optimal.  The limiting constant $4$ is established for the
explicit three-cell family and gives a lower obstruction to any universal
small-defect theory, but the best universal coefficient remains open.

The stability mechanism is also more specific than perturbation to a nearby
exact configuration.  Neither decomposition is altered.  Instead, a
near-extremal value of a coarsening gap forces the already-existing
cellwise gap profile to concentrate on two labels.  This is why the result is
naturally an inverse theorem for an inequality rather than an
almost-commuting approximation theorem.

\section{Relation to existing operator and matrix theory}
\label{sec:related}

The geometric formulation belongs first to the classical theory of relative
positions of subspaces.  Principal angles and principal vectors go back much
further than the present problem; computational and projector-based
treatments include Bj{\"o}rck--Golub, Miao--Ben-Israel, and
Gal{\'a}ntai--Heged{\H{u}}s
\cite{BjorckGolub1973,MiaoBenIsrael1992,GalantaiHegedus2006}.
For equal-dimensional subspaces, the operator norm of the difference of
their orthogonal projections is the sine of the largest principal angle.
Gal{\'a}ntai's treatment of pairs of orthogonal projections gives a
particularly relevant account of this geometry and its relation to the
Halmos--Wedin canonical forms \cite{Galantai2008}; see also
Halmos' two-subspaces theorem and the survey of B{\"o}ttcher--Spitkovsky
\cite{Halmos1969TwoSubspaces,BottcherSpitkovsky2010}.

The present problem differs from the classical two-subspace setting because
its datum is a pair of complete labelled decompositions and the observable is
not merely a single gap.  It compares the collection of cell gaps
\[
  \bigl(\|P_i-Q_i\|_{\op}\bigr)_{i\in I}
\]
with the maximal gap over all matched coarsenings
\[
  S\longmapsto\|P_S-Q_S\|_{\op}.
\]
The equality and near-equality questions therefore concern how a
finite family of subspace discrepancies can coherently aggregate.

In the equivalent commutator formulation there is extensive work on pairs of
projections.  Li studied operators representable as commutators of orthogonal
projections, and later work treats related representation and closure
questions
\cite{Li2004,ShiJi2019,MarcouxRadjaviZhang2024}.
Andruchow and Di Iorio y Lucero study geometric classes of projections whose
commutator with a fixed projection belongs to a Schatten ideal
\cite{AndruchowDiIorio2021}.  Recent work of Chaile and Chiumiento studies the
geometry of pairs of projections with a fixed commutator
\cite{ChaileChiumiento2026}.  These questions concern the geometry or
representation of a pair of projections, rather than a finite labelled
decomposition and concentration of its gap profile.

The numerical value $1/2$ also appears in the classical estimate
$\|[P,Q]\|_{\op}\le1/2$ for two orthogonal projections.  This is a different
inequality: Theorem~\ref{thm:envelope} compares a maximal coarsening gap with
the adaptive quadratic budget
$\sum_i\|P_i-Q_i\|_{\op}^2$.  The envelope is elementary once the two-sided
cut estimate has been isolated; the nontrivial part of the present work is
the inverse theory of its equality and near-equality regimes.

Sharp commutator inequalities for arbitrary matrices provide another
methodological comparison.  The B{\"o}ttcher--Wenzel inequality is a
Frobenius-norm commutator inequality with a substantial extremal theory
\cite{BottcherWenzel2008}.  The norms and objects here are different, but the
general paradigm of classifying equality and near equality is similar.

Finally, the conclusion of Theorem~\ref{thm:stability} is distinct from
almost-commuting approximation results.  Lin's theorem and quantitative work
of Hastings ask whether matrices with small commutator can be approximated
by exactly commuting matrices
\cite{Lin1997AlmostCommuting,FriisRordam1996,Hastings2009}.  For unitaries,
topological obstructions and effective variants lead to a different theory
\cite{Voiculescu1983,ExelLoring1989,ExelLoring1991,
DorOnHallKachkovskiy2025,Herrera2024,Herrera2026}.
Here the projections and the unitary are left fixed; near equality instead
forces concentration of the existing finite defect profile.

To our knowledge, the specific finite-family statement proved here---maximal
matched-coarsening gap versus the quadratic cell-gap budget, together with
two-cell equality rigidity, a cardinality-free inverse theorem, and a
three-cell inverse-square scale obstruction---does not follow from the
results above.  We make no claim that the elementary envelope inequality
itself has not appeared in equivalent block-operator language.

\section{Conclusion}

For two labelled finite orthogonal decompositions with matching cell
dimensions, let $p_i=\|P_i-Q_i\|_{\op}^2$ and
$\dd^2=\sum_i p_i$.  The maximal gap between corresponding coarsenings
satisfies
\[
  \gmax^2\le\frac{\dd^2}{2}.
\]
At positive defect, equality is rigid: exactly two cells carry nonzero gap,
and each carries half of the total quadratic budget.

Near equality has a dimension- and cardinality-free quantitative form.  If
\[
  \gmax^2\ge(1-\eta)\frac{\dd^2}{2},
  \qquad 0\le\eta\le\frac15,
\]
then the optimal relative tail outside two cells satisfies
\[
  \tau\le
  2\eta+\frac{36\eta}{(1-3\eta)\dd^2}.
\]
A one-parameter three-cell family shows that the exponent of $\dd$ is
unavoidable: as the two decompositions coalesce,
\[
  \dd_t^2\frac{\tau_t}{\eta_t}\longrightarrow4.
\]
Thus the equality classification, the inverse stability theorem, and the
scale obstruction form a single near-extremal theory for gaps of finite
orthogonal decompositions.

\section*{Declaration of generative AI and AI-assisted technologies in the manuscript preparation process}

During the preparation of this work, the author used OpenAI ChatGPT to assist
with organization of the manuscript, language editing, literature
organization, and consistency checks between the mathematical exposition and
the Lean formalization. After using this service, the author reviewed and
edited the content as needed and takes full responsibility for the content of
the publication.

\bibliographystyle{plain}
\bibliography{references_v04}

\end{document}
